%
%
% report_2025.tex
%
%
%
\documentclass[fontsize=11pt]{article} % Note: Changed from jlreq to article for English typesetting, but you can revert to jlreq if you prefer.
\usepackage{subfiles}
\usepackage{amsmath, amssymb, braket} % AMS math
\usepackage{lmodern} % Font for math
\usepackage{graphicx,color}
\usepackage{subcaption} % Captions for figures
\usepackage{bm} % Bold math
\usepackage{quantikz} % Quantum circuit diagrams
\usepackage[version=4]{mhchem} % Chemical formulas
\usepackage[backend=biber, style=numeric, maxnames=5, doi=false, url=true, eprint=true]{biblatex}
\addbibresource{quantum.bib}
\makeatletter
\let\mkbibemph=\mkbibitalic
\let\blx@imc@mkbibemph=\blx@imc@mkbibitalic
\makeatother
\setlength{\biblabelsep}{0.5em}

%\usepackage{pdfx}  % side-effect: configures hyperref
\usepackage{hyperref} % Link embedding, PDF table of contents generation
\hypersetup{
    colorlinks=true,
    linkcolor=blue,
    urlcolor=cyan,
    pdftitle={Time Evolution Simulation of 1D and 2D Hydrogen Atom Models Using Full Quantum Circuits: Error Analysis and Parameter Optimization},
    pdfauthor={Hideo Takahashi},
    pdfsubject={Research Paper},
    pdfkeywords={Quantum Computer, First Quantization, ab initio, Simulator, Coulomb Potential, Quantum Chemistry},
    bookmarksnumbered=true,
    bookmarksopen=true
}
\usepackage{xr}
\externaldocument{report_2025_sm}
\usepackage{paper2025}

\begin{document}
\title
{Time Evolution Simulation of a Hydrogen Atom Using Full Quantum Circuits}

\author{Hideo Takahashi, Tatsuya Tomaru, Toshiyuki Hirano, Saisei Tahara, Fumitoshi Sato}
\date{\today}
\maketitle

\begin{abstract}
In this study, we constructed a quantum circuit for a first-quantization-based
Hamiltonian simulator capable of operating on a small-scale quantum computer
emulator with fewer than 30 qubits.

To construct the circuit within limited quantum resources (number of qubits), we
applied Uniformly Controlled Rotation (UCR) circuits to calculate the Coulomb
potential. Furthermore, as a guideline for obtaining highly accurate
computational results, we established a method for handling the singularity
(pole) of the Coulomb potential and optimizing the discretization parameters.
Regarding the singularity of the Coulomb potential, we demonstrated that
divergence can be avoided and errors minimized by determining an effective value
from a simple geometric average of the potential within a minute circular region
centered on the pole.

Additionally, under the constraint of a given number of qubits, we proposed a
method to determine the optimal spatial grid spacing by evaluating the tradeoff
between the spatial discretization error caused by grid coarseness and the
domain truncation error due to a finite domain length. For the time evolution
step size, rather than relying on the formal Trotter error bound that assumes
the worst-case scenario, we clarified the conditions required to maintain
sufficient accuracy while avoiding an unnecessary increase in quantum circuit
depth. This was achieved based on the sampling theorem, focusing on the maximum
frequency component of the wave function.
\end{abstract}

\subfile{report2025_1_intro.tex}
\subfile{report2025_2_method.tex}
\subfile{report2025_3_results.tex}
\subfile{report2025_4_analysis.tex}

\section{Conclusion}
In this study, we presented a circuit configuration capable of simulating the
full time evolution of a two-dimensional hydrogen atom model on a Qiskit
emulator. To the best of our knowledge, this is the first report to implement
grid-based time evolution calculations, including the Coulomb potential,
entirely as a quantum circuit and to confirm its operation on an emulator.

Regarding the handling of the Coulomb potential singularity, we proposed a
method using $\reff=\deltar/4$ as the effective distance to determine the
effective potential at the grid point containing the singularity in the 2D
model. Numerical experiments confirmed not only that this value minimizes the
energy error of the system, but also that it coincides with the value derived
from the geometric area average around the singularity, clarifying the physical
validity of this approach in a discretized space.

Furthermore, we established a procedure for determining simulation parameters
that minimize errors under a given number of qubits ($\nb$). We demonstrated
that the spatial grid spacing $\deltar$ can be determined as the optimal value
and the dynamic range of the simulation by balancing the tradeoff between
spatial discretization error and domain truncation error. 

Moreover, regarding the time step $\deltat$ of the Trotter decomposition, we
numerically confirmed that it can be determined not from the formal upper bound
of the Trotter error assuming the worst case, but from the requirement of the
sampling theorem (half-period condition) based on the maximum frequency
component of the wave function. Particularly when targeting states near the
ground state, we demonstrated that within the range satisfying this condition,
the spatial error becomes dominant over the time error (i.e., the error
saturates), providing a guideline to avoid unnecessary increases in quantum
circuit depth.

These findings are expected to serve as highly practical and powerful guidelines
for designing and executing quantum simulations using the grid method under
limited quantum resources.

\printbibliography[heading=bibintoc, title={References}]

\subfile{report2025_a_appendix.tex}

\end{document}
